% 赣南科技学院毕业论文LaTeX模板
\documentclass[a4paper,12pt,openany,twoside]{ctexbook}

% 页面设置
% 要求：上下左右边距2.5/2.5/2.5/2.0cm，页眉1.5cm，页脚1.75cm
% Word设置：页眉距纸张顶部1.5cm，正文距纸张顶部2.5cm
% asymmetric: 确保奇偶页使用相同的左右边距（不交换）
\usepackage[
    top=2.5cm,           % 上边距2.5cm（正文距纸张顶部）
    bottom=2.5cm,        % 下边距2.5cm
    left=2.5cm,          % 左边距2.5cm（所有页面）
    right=2.0cm,         % 右边距2.0cm（所有页面）
    headheight=0.5cm,    % 页眉高度（五号字约14bp≈0.5cm）
    headsep=0.5cm,       % 页眉到正文的距离
    footskip=0.75cm,     % 正文到页脚的距离
    asymmetric           % 奇偶页使用相同左右边距
]{geometry}
% 说明：
% - 页眉距纸张顶部 = top - headheight - headsep = 2.5 - 0.5 - 0.5 = 1.5cm ✓
% - 页脚距纸张底部 = bottom - footskip = 2.5 - 0.75 = 1.75cm ✓
% - 设置"奇偶页不同"通过twoside选项实现，asymmetric确保边距不交换

% 字体设置
\usepackage{fontspec}
\setmainfont{Times New Roman}
\setCJKmainfont{SimSun}[AutoFakeBold=2.5,ItalicFont=KaiTi]
\setCJKsansfont{SimHei}[AutoFakeBold=2.5]
\setCJKmonofont{FangSong}

% 其他必要的包
\usepackage{graphicx}      % 图片
\usepackage{amsmath}       % 数学公式
\usepackage{amssymb}       % 数学符号
\usepackage{booktabs}      % 表格
\usepackage{longtable}     % 长表格
\usepackage{multirow}      % 表格合并
\usepackage{caption}       % 图表标题
\usepackage{subcaption}    % 子图
\usepackage{enumitem}      % 列表
\usepackage{fancyhdr}      % 页眉页脚
\usepackage{titlesec}      % 标题格式
\usepackage{titletoc}      % 目录格式
\usepackage[hidelinks]{hyperref}  % 超链接
\usepackage{xcolor}        % 颜色
\usepackage{listings}      % 代码
\usepackage{setspace}      % 行距
\usepackage{indentfirst}   % 首行缩进

% 设置中文段落缩进
\setlength{\parindent}{2em}

% 章节标题格式设置
\ctexset{
    chapter={
        format={\centering\heiti\mdseries\zihao{-3}},  % 小三号黑体居中，不加粗
        nameformat={},
        titleformat={},
        number=\chinese{chapter},
        name={第,章},
        beforeskip=1\baselineskip,     % 段前1行       
        afterskip=1\baselineskip,    % 段后1行
        fixskip=true
    },
    % contentsname 在目录部分手动设置
    listfigurename={插图目录},
    listtablename={表格目录},
    section={
        format={\songti\bfseries\zihao{4}},  % 四号宋体加粗，顶格，段前段后0.5行
        beforeskip=0.5\baselineskip,
        afterskip=0.5\baselineskip,
        indent=0pt,       % 顶格
        aftername=\       % 序号后空一格
    },
    subsection={
        format={\kaishu\bfseries\zihao{4}},  % 四号楷体加粗，顶格
        beforeskip=0.5\baselineskip,
        afterskip=0.5\baselineskip,
        indent=0pt,       % 顶格
        aftername=\       % 序号后空一格
    }
}

% 图表标题设置（五号黑体）
\DeclareCaptionFont{heiti}{\heiti}
\DeclareCaptionFont{wuhao}{\zihao{5}}
\captionsetup{
    font={heiti,wuhao},  % 五号黑体
    labelsep=space,
    skip=5pt
}
\captionsetup[figure]{
    position=bottom,
    aboveskip=5pt,
    belowskip=0pt
}
\captionsetup[table]{
    position=top,
    aboveskip=0pt,
    belowskip=5pt
}

% 重新定义图表编号格式
\renewcommand{\thefigure}{\thechapter-\arabic{figure}}
\renewcommand{\thetable}{\thechapter-\arabic{table}}
\renewcommand{\theequation}{\thechapter-\arabic{equation}}

% 目录格式设置
% - 各章节标题：四号字体，不加粗，1.5倍行距
% - 第一级标题（章）：四号黑体，缩进2个全角空格（相对于左边距2.5cm）
% - 第二级标题（节）：四号宋体，缩进4个全角空格
% - 第三级标题（小节）：四号楷体，缩进6个全角空格
% - 题目和页码间用"…………"相连
% 设置目录右边距，确保遵循全局右边距2.0cm
\contentsmargin{0pt}
% 章：左缩进6em（标签起始2em），标签宽4em
\titlecontents{chapter}[6em]{\heiti\zihao{4}\vspace{6pt}}
    {\contentslabel[\thecontentslabel\hfill]{4em}}{\hspace{-4em}}{\titlerule*[0.5pc]{.}\contentspage}
% 节：左缩进6.5em（标签起始4em），标签宽2.5em
\titlecontents{section}[6.5em]{\songti\zihao{4}}
    {\contentslabel{2.5em}}{}{\titlerule*[0.5pc]{.}\contentspage}
% 小节：左缩进9.5em（标签起始6em），标签宽3.5em
\titlecontents{subsection}[9.5em]{\kaishu\zihao{4}}
    {\contentslabel{3.5em}}{}{\titlerule*[0.5pc]{.}\contentspage}

% 页眉页脚设置
\pagestyle{fancy}
\fancyhf{}
\fancyhead[CO]{\songti\zihao{5}赣南科技学院2025届毕业论文}  % 奇数页页眉
\fancyhead[CE]{\songti\zihao{5}学生姓名：论文题目}  % 偶数页页眉
\fancyfoot[C]{\zihao{5}\thepage}  % 页码居中
\renewcommand{\headrulewidth}{0.4pt}
\renewcommand{\footrulewidth}{0pt}

% 前置部分页面样式（摘要、目录，使用罗马数字页码）
\fancypagestyle{frontmatterstyle}{
    \fancyhf{}
    \fancyhead[CO]{\songti\zihao{5}赣南科技学院2025届毕业论文}
    \fancyhead[CE]{\songti\zihao{5}学生姓名：论文题目}
    \fancyfoot[C]{\zihao{5}\Roman{page}}  % 罗马数字页码
    \renewcommand{\headrulewidth}{0.4pt}
    \renewcommand{\footrulewidth}{0pt}
}

% 重新定义plain样式（用于章节首页），使其显示页眉
\fancypagestyle{plain}{
    \fancyhf{}
    \fancyhead[CO]{\songti\zihao{5}赣南科技学院2025届毕业论文}
    \fancyhead[CE]{\songti\zihao{5}学生姓名：论文题目}
    \fancyfoot[C]{\zihao{5}\thepage}
    \renewcommand{\headrulewidth}{0.4pt}
    \renewcommand{\footrulewidth}{0pt}
}

% 参考文献设置
% 使用手动 thebibliography 环境时，不需要 gbt7714 包
% \usepackage[sort&compress]{gbt7714}
% \bibliographystyle{gbt7714-numerical}

% 代码列表设置
\lstset{
    basicstyle=\ttfamily\small,
    keywordstyle=\color{blue},
    commentstyle=\color{gray},
    stringstyle=\color{red},
    numbers=left,
    numberstyle=\tiny,
    stepnumber=1,
    numbersep=8pt,
    breaklines=true,
    frame=single,
    xleftmargin=2em,
    xrightmargin=2em,
    tabsize=4
}

% ==================== 论文信息 ====================
\newcommand{\thesistitle}{基于Spring Boot的图书管理系统的设计与实现}
\newcommand{\thesissubtitle}{} % 副标题（如果有）
\newcommand{\thesisauthor}{学生姓名}
\newcommand{\studentid}{学号}
\newcommand{\supervisor}{指导教师}
\newcommand{\college}{学院}
\newcommand{\major}{专业}
\newcommand{\classno}{班级}
\newcommand{\thesisyear}{2025}

% 英文信息
\newcommand{\thesistitleEN}{The Design and Implementation of a Library Management System Based on Spring Boot}
\newcommand{\thesissubtitleEN}{}
\newcommand{\thesisauthorEN}{Student Name}
\newcommand{\supervisorEN}{Supervisor Name}
\newcommand{\collegeEN}{College Name}
\newcommand{\majorEN}{Major Name}

% ==================== 文档开始 ====================
\begin{document}

% 前置部分（摘要、目录）使用罗马数字页码
\frontmatter
\pagestyle{frontmatterstyle}
\setcounter{page}{1}

% ==================== 中文摘要 ====================
% 格式要求：
% - 论文题目：小一号黑体居中，题目较长可分多行书写，单倍行距，与下面内容不空行
% - "摘  要"：三号黑体居中，"摘"与"要"之间空两格，与下面内容单倍行距
% - 摘要内容：小四号宋体，20磅行距，首行缩进2字符，400字以内
% - 与下面内容单倍行距
% - "关键词："：小四号黑体，顶格
% - 关键词内容：小四号宋体，不加粗，单倍行距，中文分号隔开，最后无标点
\chapter*{}

\thispagestyle{frontmatterstyle}
\begin{center}
    {\heiti\zihao{-1}\thesistitle}
    \par
    {\heiti\zihao{3}摘~~要}
\end{center}

{\songti\zihao{-4}\setlength{\baselineskip}{20pt}
\setlength{\parindent}{2em}
\indent 这里是中文摘要内容。摘要一般在400字以内，语言力求精炼。摘要应概括地反映出本论文的主要内容，主要说明本论文的研究目的、内容、方法、成果和结论。要突出本论文的创新之处，语言力求精炼。为了便于文献检索，应在摘要下方另起一行注明论文的关键词。
}

\vspace{\baselineskip}
\noindent{\heiti\zihao{-4}关键词：}{\songti\zihao{-4}关键词1；关键词2；关键词3；关键词4；关键词5}

\clearpage

% ==================== 英文摘要 ====================
% 格式要求：
% - 英文论文题目：小一号Times New Roman加粗居中，题目较长可分多行书写，与下面内容单倍行距
%   实词首字母大写，冠词/短介词/并列连词首字母小写，第一个单词首字母大写
% - "ABSTRACT"：三号Times New Roman加粗居中，与下面内容单倍行距
% - 摘要内容：小四号Times New Roman，20磅行距，首行缩进4个半角空格
% - 空1行
% - "Keywords:"：小四号Times New Roman加粗，顶格
% - 关键词内容：小四号Times New Roman，不加粗，英文分号隔开
\chapter*{}

\thispagestyle{frontmatterstyle}
\begin{center}
    {\bfseries\zihao{-1}\thesistitleEN}
    \par
    {\bfseries\zihao{3}ABSTRACT}
\end{center}

{\zihao{-4}\setlength{\baselineskip}{20pt}
\hspace{4ex}This is the English abstract. The content should correspond to the Chinese abstract. The abstract should be concise and generally within 400 words.
}

\vspace{\baselineskip}
\noindent{\bfseries\zihao{-4}Keywords:} {\zihao{-4}Keyword1; Keyword2; Keyword3; Keyword4; Keyword5}

\clearpage

% ==================== 目录 ====================
% 格式要求：
% - "目录"：三号黑体居中，"目"与"录"之间两个半角空格，与下面内容空一行
% - 目录内容：各章节标题，不超过三级，四号字体，不加粗，1.5倍行距
\pagestyle{frontmatterstyle}
\chapter*{}

\thispagestyle{frontmatterstyle}
\begin{center}
    {\heiti\zihao{3}目~~录}
\end{center}
\vspace{0.5\baselineskip}  % 与下面内容空一行
{
\makeatletter
\@starttoc{toc}
\makeatother
}
\clearpage

% ==================== 正文 ====================
\mainmatter
\pagestyle{fancy}  % 切换回正常页面样式
\setcounter{page}{1}  % 正文页码从1开始（阿拉伯数字）

% 设置正文统一格式：中文小四号宋体，英文小四号Times New Roman，行距固定值20磅，标准字符间距，首行缩进2字符，段前段后间距均为0行
\songti\zihao{-4}
\setlength{\baselineskip}{20pt}
\setlength{\parindent}{2em}
\setlength{\parskip}{0pt}

% 修改偶数页页眉为实际信息
\fancyhead[CE]{\songti\zihao{5}\thesisauthor：\thesistitle}

% 重新定义正文部分的plain样式（章节首页），使其显示正确的页眉
\fancypagestyle{plain}{
    \fancyhf{}
    \fancyhead[CO]{\songti\zihao{5}赣南科技学院2025届毕业论文}
    \fancyhead[CE]{\songti\zihao{5}\thesisauthor：\thesistitle}
    \fancyfoot[C]{\zihao{5}\thepage}
    \renewcommand{\headrulewidth}{0.4pt}
    \renewcommand{\footrulewidth}{0pt}
}

\chapter{绪论}

% 设置正文格式：小四号宋体，20磅行距，首行缩进2字符，段前段后0行

这里是正文内容。正文中中文采用小四号宋体，英文采用小四号Times New Roman字体，行距采用固定值20磅，标准字符间距，首行缩进2字符。

\section{课题背景和意义}

这是二级标题的内容（四号宋体加粗，顶格，段前段后间距0.5行）。内容采用小四号宋体，20磅行距，首行缩进2字符。

\subsection{国内研究现状}

这是三级标题的内容（四号楷体加粗，顶格，段前段后间距0.5行）。内容采用小四号宋体，20磅行距，首行缩进2字符。

\section{论文的组织结构}

本论文共分为五章，各章节内容安排如下：

第一章为绪论，介绍了课题的研究背景、意义以及国内外研究现状。

第二章介绍了本系统所使用的相关技术。

第三章详细阐述了系统的需求分析和总体设计方案。

第四章介绍了系统的具体实现过程。

第五章对全文进行总结，并对未来的工作进行展望。

\section{数字编号示例}

正文中的数字编号示例（数字及符号使用Times New Roman体）：

\begin{enumerate}[label=\arabic*.,leftmargin=2em,itemsep=0pt]
    \item 一级编号，数字加小数点"."加一个空格
    \item 总体思路说明
\end{enumerate}

\begin{enumerate}[label=(\arabic*),leftmargin=2.5em,itemsep=0pt]
    \item 二级编号，英文下的括号(数字)加一个空格
    \item 工作要求说明
\end{enumerate}

\begin{enumerate}[label=\arabic*),leftmargin=2.5em,itemsep=0pt]
    \item 三级编号，数字后跟英文下的括号)再加一个空格
    \item 集合地点说明
\end{enumerate}

\begin{enumerate}[label={\heiti ①},leftmargin=2em,itemsep=0pt]
    \item 四级编号，①加一个空格
    \item 支持串口虚拟功能
\end{enumerate}

\begin{enumerate}[label=\Alph*.,leftmargin=2em,itemsep=0pt]
    \item 五级编号，大写字母加小数点"."加一个空格
    \item 全面总结
\end{enumerate}

\begin{enumerate}[label=\alph*.,leftmargin=2em,itemsep=0pt]
    \item 六级编号，小写字母加小数点"."加一个空格
    \item 开源技术
\end{enumerate}

\section{图表公式示例}

\subsection{图片示例}

图片应放置在合适位置，图序及图名居中置于图的下方。图序编号按章内连续编号，格式为"图M-N"，其中M为章号，N为图序号。图注采用五号黑体，居中。

\begin{figure}[htbp]
    \centering
    % 将图片放在figures文件夹中，然后取消下面的注释
    % \includegraphics[width=0.6\textwidth]{figures/example.png}
    
    % 临时占位框
    \fbox{\parbox{0.6\textwidth}{\centering\vspace{4cm}此处为图片位置\\请将图片放入figures文件夹}}
    \caption{瀑布模型示意图}
    \label{fig:waterfall}
\end{figure}

引用图片示例：如图\ref{fig:waterfall}所示，瀑布模型是一种经典的软件开发模型。

图与上方正文之间不空行，图与下方正文空一行。

\vspace{\baselineskip}

\subsection{表格示例}

表序及表名置于表的上方，居中。表序编号按章内连续编号，格式为"表M-N"。表序、表名和表内内容采用五号黑体字。

\begin{table}[htbp]
    \centering
    \caption{换热器不同级别损失表}
    \label{tab:heat-exchanger}
    \begin{tabular}{ccc}
        \toprule
        \textbf{换热器} & \textbf{热边压降损失（千卡）} & \textbf{冷边压降损失（千卡）} \\
        \midrule
        初级 & 2974.37 & 2931.52 \\
        次级 & 2924.65 & 3798.76 \\
        \bottomrule
    \end{tabular}
\end{table}

引用表格示例：如表\ref{tab:heat-exchanger}所示。表与下方正文空一行。

\vspace{\baselineskip}

\subsection{公式示例}

公式编号用括号括起写在右边行末，格式为(M-N)，其中M为章号，N为公式序号。公式中的英文字母要有一行的间距，公式与正文之间不需空行。

质能方程：
\begin{equation}
    E = mc^2
    \label{eq:mass-energy}
\end{equation}

热力学第一定律：
\begin{equation}
    \Delta U = Q - W
    \label{eq:thermodynamics}
\end{equation}

引用公式示例：如式(\ref{eq:mass-energy})所示，爱因斯坦的质能方程揭示了质量与能量的关系。

行内公式示例：根据公式 $E=mc^2$ 可知，质量和能量是可以相互转换的。

\chapter{系统技术论述}

本章主要介绍系统开发所使用的相关技术，包括开发框架、数据库技术、前端技术等。这些技术的合理运用为系统的高效开发和稳定运行提供了保障。

\section{开发框架介绍}

本系统采用Spring Boot框架进行开发。Spring Boot是Spring家族中的一个全新框架，其设计目的是用来简化Spring应用的初始搭建以及开发过程。该框架使用了特定的方式来进行配置，从而使开发人员不再需要定义样板化的配置。

\subsection{Spring Boot特点概述}

Spring Boot具有以下主要特点：

\begin{enumerate}[label=(\arabic*),leftmargin=2.5em,itemsep=0pt]
    \item 创建独立的Spring应用程序
    \item 嵌入的Tomcat，无需部署WAR文件
    \item 简化Maven配置
    \item 自动配置Spring
    \item 提供生产就绪型功能
    \item 绝对没有代码生成和对XML没有要求配置
\end{enumerate}

\section{数据库技术}

本系统采用MySQL数据库进行数据存储和管理。MySQL是一个关系型数据库管理系统，由瑞典MySQL AB公司开发，目前属于Oracle旗下产品。

\section{前端技术}

前端采用Vue.js框架开发，Vue.js是一套构建用户界面的渐进式框架，易于上手，便于与第三方库或既有项目整合。

\chapter{系统需求分析与设计}

本章对系统进行需求分析，明确系统的功能需求和非功能需求，在此基础上进行系统的总体设计和详细设计。

\section{需求分析}

需求分析是系统开发的基础，通过对用户需求的深入分析，可以明确系统应该具备哪些功能，为后续的设计和实现指明方向。

\section{系统总体设计}

根据需求分析的结果，对系统进行总体架构设计，包括系统的层次结构、模块划分、数据流向等。

\section{数据库设计}

数据库设计是系统设计的重要组成部分，合理的数据库设计可以提高系统的性能和可维护性。

\chapter{系统实现}

本章详细介绍系统各个功能模块的具体实现过程，包括核心代码和关键算法的实现。

\section{用户管理模块实现}

用户管理模块负责系统用户的注册、登录、权限管理等功能。

\section{数据管理模块实现}

数据管理模块实现了对业务数据的增删改查等基本操作。

\chapter{系统测试}

系统测试是保证系统质量的重要环节。本章介绍了系统的测试方案、测试用例以及测试结果。

\section{测试环境}

本系统的测试环境配置如下：操作系统为Windows 10，数据库为MySQL 8.0，浏览器为Chrome最新版本。

\section{功能测试}

对系统的各个功能模块进行了详细的功能测试，测试结果表明系统各项功能均能正常运行。

\chapter{总结与展望}

本文设计并实现了一个基于Spring Boot的图书管理系统。该系统实现了图书信息管理、用户管理、借阅管理等功能，界面友好，操作简便，能够满足中小型图书馆的日常管理需求。

通过本次毕业设计，我掌握了完整的软件开发流程，从需求分析、系统设计、编码实现到测试部署，每个环节都得到了充分的锻炼。同时也深入学习了Spring Boot、Vue.js等主流开发技术，为今后的学习和工作打下了坚实的基础。

由于时间和个人能力有限，系统还存在一些不足之处，有待进一步完善：

\begin{enumerate}[label=(\arabic*),leftmargin=2.5em,itemsep=0pt]
    \item 系统的性能优化还有提升空间
    \item 可以增加更多的统计分析功能
    \item 移动端适配有待加强
\end{enumerate}

在今后的工作中，我将继续完善系统功能，提升系统性能，使其更加完善和实用。


% ==================== 参考文献 ====================
\clearpage
\phantomsection
\addcontentsline{toc}{chapter}{参考文献}

\begin{center}
    {\heiti\zihao{3}参考文献}
\end{center}

\vspace{\baselineskip}  % 空一行

% 参考文献内容：小四号宋体，20磅行距
{\songti\zihao{-4}\setlength{\baselineskip}{20pt}
\begin{thebibliography}{99}

\bibitem{ref1}
何海浪. 机器人铣削加工轨迹规划与颤振稳定性研究[J]. 计算机应用, 2020, 7(3): 167-173, 186.

\bibitem{ref2}
严蔚敏. 数据结构（第四版）[M]. 北京: 清华大学出版社, 2019: 234-236.

\bibitem{ref3}
张三, 李四, 王五. 基于Spring Boot的图书管理系统设计与实现[J]. 计算机工程与应用, 2021, 57(12): 245-250.

\bibitem{ref4}
Smith J, Johnson M, Williams R. Design and implementation of library management system[J]. Computer Science, 2022, 15(2): 123-135.

\bibitem{ref5}
赵六. Spring Boot实战[M]. 北京: 机械工业出版社, 2020: 100-150.

\bibitem{ref6}
Brown A, Davis B. Vue.js in Action[M]. New York: Manning Publications, 2021: 200-250.

\bibitem{ref7}
孙七. 基于微服务架构的图书管理系统研究[D]. 北京: 北京大学, 2021: 30-45.

\bibitem{ref8}
Wilson C D. Research on Database Optimization[D]. Boston: MIT, 2020: 50-70.

\bibitem{ref9}
周八. 软件系统架构设计方法研究[C]. 北京: 电子工业出版社, 2022: 15-20.

\bibitem{ref10}
Miller E F, Anderson G H. Software Testing Methods[C]. London: IEEE Press, 2021: 80-90.

\bibitem{ref11}
国家标准化管理委员会. GB/T 7714-2015 信息与文献参考文献著录规则[S]. 北京: 中国标准出版社, 2015.

\bibitem{ref12}
陈十一. MySQL数据库优化技术研究[J]. 软件学报, 2022, 33(4): 1234-1241.

\bibitem{ref13}
Thompson R, White S, Black T, et al. Modern web application development[J]. ACM Computing Surveys, 2022, 54(3): 456-480.

\bibitem{ref14}
刘十二. Java Web开发实战[M]. 上海: 上海科技出版社, 2021: 78-89.

\bibitem{ref15}
Green P, Blue Q. RESTful API Design[M]. Cambridge: O'Reilly Media, 2020: 150-170.

\end{thebibliography}

\textbf{注意：}参考文献一般不应少于15篇，一般应为近5年的文献。除个别专业外，一般应有5篇以上外文参考文献。作者只写到第三位，余者写"等"，英文作者超过3人写"et al"（斜体）。在正文中要有引用标注，如×××[1]。


% ==================== 附录 ====================
\clearpage
\phantomsection
\addcontentsline{toc}{chapter}{附录}

\begin{center}
    {\heiti\zihao{3}附~~录}
\end{center}

\vspace{\baselineskip}  % 空一行

% 附录内容：字体段落格式与正文要求一致
附录的有无根据说明书（设计）情况而定，内容一般包括正文内不便列出的冗长公式推导、符号说明（含缩写）、计算机程序等，字体段落格式与正文要求一致。附录中有程序源代码的因篇幅限制可酌情考虑内容的序号。

\textbf{需要特别强调：}如无确实必要，不要加附录，论文重复率检测是会检测附录的。

\section*{附录A 核心代码}

以下是系统的核心代码片段：

\begin{lstlisting}[language=Java, caption=用户控制器代码]
@RestController
@RequestMapping("/api/user")
public class UserController {
    
    @Autowired
    private UserService userService;
    
    @PostMapping("/login")
    public Result login(@RequestBody LoginDTO loginDTO) {
        User user = userService.login(loginDTO);
        if (user != null) {
            return Result.success(user);
        }
        return Result.error("用户名或密码错误");
    }
    
    @GetMapping("/list")
    public Result getUserList() {
        List<User> users = userService.getAllUsers();
        return Result.success(users);
    }
}
\end{lstlisting}

\section*{附录B 数据库表结构}

以下是系统主要数据表的结构说明：

\textbf{用户表（user）：}存储用户的基本信息，包括用户ID、用户名、密码、邮箱、创建时间等字段。

\textbf{图书表（book）：}存储图书的详细信息，包括图书ID、书名、作者、出版社、ISBN、分类、库存等字段。

\textbf{借阅记录表（borrow\_record）：}存储借阅记录信息，包括记录ID、用户ID、图书ID、借阅时间、归还时间、状态等字段。

\section*{附录C 缩写说明}

\begin{itemize}[leftmargin=2em,itemsep=0pt]
    \item API: Application Programming Interface，应用程序接口
    \item DTO: Data Transfer Object，数据传输对象
    \item MVC: Model-View-Controller，模型-视图-控制器
    \item ORM: Object-Relational Mapping，对象关系映射
    \item REST: Representational State Transfer，表现层状态转换
    \item SQL: Structured Query Language，结构化查询语言
\end{itemize}


% ==================== 致谢 ====================
\clearpage
\phantomsection
\addcontentsline{toc}{chapter}{致谢}

\begin{center}
    {\heiti\zihao{3}致~~谢}
\end{center}

\vspace{\baselineskip}  % 与内容空一行

% 致谢内容：小四号宋体，20磅行距，首行缩进2字符，限1页
时光荏苒，四年的大学生活即将结束。回首这四年的学习历程，我收获颇丰，这离不开各位老师、同学和家人的帮助与支持。

在此，我要特别感谢我的指导教师李红旗教授和张爱国老师。在整个毕业设计过程中，从选题、开题、研究到论文撰写的各个阶段，李老师和张老师都给予了我悉心的指导和热情的帮助。他们严谨的治学态度、渊博的专业知识和科学的工作方法使我受益匪浅，对我今后的学习和工作产生了深远的影响。

感谢赣南科技学院电子信息工程学院的各位老师，是你们的辛勤付出和谆谆教诲，让我在专业知识和实践能力方面得到了全面的提升。

感谢我的同学们，在这四年的学习生活中，我们相互帮助、共同进步，结下了深厚的友谊。特别感谢我的室友们，你们的陪伴让我的大学生活更加充实和快乐。

感谢我的父母和家人，是你们无私的爱和无条件的支持，让我能够安心学习、顺利完成学业。你们的理解和鼓励是我前进的动力。

最后，感谢所有关心和帮助过我的人！


\end{document}

